% Part: methods
% Chapter: proofs
% Section: inference-patterns

\documentclass[../../../include/open-logic-section]{subfiles}

\begin{document}

\olfileid{mth}{prf}{pat}

% BN-SRC-833: inference steps are distinct from temporary assumptions.
\olsection{অনুমিতির রীতি}

একটি প্রমাণ অনেকগুলি স্বতন্ত্র অনুমিতি-ধাপ দিয়ে গঠিত।
কোনো সিদ্ধান্ত টানার সময় সাধারণত ``তাই'', ``সুতরাং'' বা
``অতএব'' লিখে তা বোঝাই। সিদ্ধান্তটি প্রমাণে আগে পাওয়া
এক বা দুটি তথ্যের উপর নির্ভর করে—সেগুলি হয়তো আগে
অনুমান করেছি, অথবা আগের কোনো অনুমিতি-ধাপে পেয়েছি।
স্পষ্টতার জন্য সেগুলিকে চিহ্নিত করতে পারি এবং নতুন ধাপে
কোন বক্তব্যগুলি ব্যবহার করছি, তা বলতে পারি। আগের কথা
থেকে নতুন সিদ্ধান্তটি \emph{কেন} আসে, তার ব্যাখ্যাও প্রায়ই
অনুমিতি-ধাপে থাকে। প্রমাণে বারবার ব্যবহৃত কয়েকটি সাধারণ
অনুমিতির রীতি নিচে দেখব। কিছু রীতি, যেমন আরোহের প্রমাণ,
আরও জটিল; সেগুলি পরে আলোচনা করব।

আগেই একটি অনুমিতির রীতি দেখেছি: সংজ্ঞা খুলে লেখা বা
প্রয়োগ করা। সংজ্ঞা খুললে সংজ্ঞেয় পদ-যুক্ত কথাটিকে
সংজ্ঞাদায়ী অংশের সাহায্যে আবার বলি। যেমন, ধরো প্রমাণে
ইতিমধ্যে $D = E$ প্রতিষ্ঠা করেছি—এটি (a)। তখন সেটের
$=$-সংজ্ঞা প্রয়োগ করে বলতে পারি: ``অতএব (a) ও সংজ্ঞা
থেকে, $D$-র প্রত্যেক !!{element} $E$-রও
!!a{element}, এবং উল্টোটিও সত্য।''

কিছুটা বিভ্রান্তিকর হলেও, কোনো সিদ্ধান্তের সমর্থন অনেক
সময় সিদ্ধান্তটি টানার মুহূর্তে নয়, তার আগেই লিখি। ধরো,
$D = E$ এখনো প্রমাণ করিনি, কিন্তু করতে চাই। লক্ষ্য
$D = E$-কে সংজ্ঞা খুলে আবার বলতে পারি: $D = E$
প্রমাণ করতে $D$-র প্রত্যেক !!{element} $E$-রও
!!a{element}, এবং উল্টোটিও, প্রমাণ করতে হবে। তাই
প্রমাণের রূপ হবে: (a) $D$-র প্রত্যেক !!{element}
$E$-রও !!a{element} প্রমাণ; (b) $E$-র প্রত্যেক
!!{element} $D$-রও !!a{element} প্রমাণ;
(c) অতএব (a), (b) ও $=$-র সংজ্ঞা থেকে $D = E$।
কিন্তু সচরাচর এভাবে লিখি না। বরং লিখতে পারি:
\begin{quote}
দেখাতে চাই $D = E$। $=$-র সংজ্ঞা অনুযায়ী, এর
জন্য দেখাতে হবে $D$-র প্রত্যেক !!{element} $E$-রও
!!a{element}, এবং উল্টোটিও।

(a) \dots ($D$-র প্রত্যেক !!{element} $E$-রও
!!a{element} হওয়ার প্রমাণ) \dots

(b) \dots ($E$-র প্রত্যেক !!{element} $D$-রও
!!a{element} হওয়ার প্রমাণ) \dots
\end{quote}

\subsection{সংযোজনের ব্যবহার}

সম্ভবত সবচেয়ে সরল অনুমিতির রীতি হলো সংযোজিত
বক্তব্যের একটি অংশকে সিদ্ধান্ত হিসেবে নেওয়া। অর্থাৎ
$p$ এবং~$q$ সত্য বলে ধরে নিলে বা প্রমাণ করলে,
আমরা $p$ (এবং $q$-ও) বলতে পারি। ধাপটি এত
মৌলিক যে অনেক সময় আলাদা করে উল্লেখ করি না।
যেমন, $D = E$-র সংজ্ঞা খুলে জেনেছি $D$-র প্রত্যেক
!!{element} $E$-রও !!a{element}, এবং উল্টোটিও।
এ থেকে বলতে পারি $E$-র প্রত্যেক !!{element}
$D$-রও !!a{element}—এটিই ``উল্টোটিও'' অংশ।

\subsection{সংযোজন প্রমাণ}

কখনো যে বক্তব্য প্রমাণ করতে হবে তা সংযোজনের আকারে
থাকে: ``$p$ এবং $q$ প্রমাণ করো।'' সে ক্ষেত্রে দুটি
কাজ করতে হবে: $p$ প্রমাণ, তারপর $q$ প্রমাণ।
স্পষ্টতার জন্য প্রমাণকে দুটি চিহ্নিত অংশে ভাগ করতে পারো।
প্রথম খসড়ায় কাগজের উপরে ``(1) $p$ প্রমাণ'' এবং
মাঝখানে ``(2) $q$ প্রমাণ'' লিখে নিতে পারো। (সরাসরি
সংযোজন প্রমাণ করতে না বললেও প্রমাণে তা লাগতে পারে।
যেমন, $D = E$ প্রমাণ করতে বললে $=$-র সংজ্ঞা খুলে
দেখবে, প্রমাণ করতে হবে $D$-র প্রত্যেক !!{element}
$E$-রও !!a{element} \emph{এবং} $E$-র প্রত্যেক
!!{element} $D$-রও !!a{element}।)

\subsection{বিযোজন প্রমাণ}

প্রমাণাধীন বক্তব্য বিযোজনের আকারে থাকলে—অর্থাৎ
``$p$ অথবা $q$''—বিকল্প দুটির একটি সত্য দেখালেই
চলে। তবে উপপাদ্যের পূর্বধারণা থেকে বিকল্পের একটিই
সরাসরি পাওয়া যায়, এমন ঘটনা খুব কম। অধিকাংশ ক্ষেত্রে
পূর্বধারণাই বিযোজক; অথবা নির্দিষ্ট ধরনের সব বস্তুর
দুটি ধর্মের একটি আছে দেখাচ্ছ, কিন্তু কিছু বস্তুর একটি
আর বাকিদের অন্যটি আছে। তখন ক্ষেত্রবিচারে প্রমাণ
কাজে লাগে (নিচে দেখো)।

\subsection{শর্তাধীন প্রমাণ}

অনেক উপপাদ্য শর্তাধীন আকারে থাকে: $p$ সত্য হলে
$q$-ও সত্য দেখাতে হয়। এমন প্রমাণ শুরু করা সহজ—
শর্তের পূর্বপক্ষ, অর্থাৎ এখানে $p$, অনুমান করে
সেখান থেকে উত্তরপক্ষ~$q$ প্রমাণ করো। উপপাদ্যে
``$p$ হলে $q$'' থাকলে, ``$p$ ধরি'' লিখে শুরু করবে
এবং শেষে~$q$ প্রমাণিত হবে।

শর্তাধীন বক্তব্য বিভিন্নভাবে বলা যায়। ``$p$ হলে
$q$''-র বদলে উপপাদ্যে ``$p$ কেবল যদি $q$'',
``$q$, যদি $p$'' বা ``$p$ শর্তে $q$'' থাকতে পারে।
সবগুলির অর্থ একই: $p$ ধরে তার থেকে~$q$ প্রমাণ
করতে হয়। মনে করো, দ্বিশর্তাধীন ``$p$ তখনই এবং
কেবল তখনই, যখন $q$'' আসলে দুটি শর্তাধীন বক্তব্য:
$p$ হলে $q$, আর $q$ হলে~$p$। তাই দুটি
শর্তাধীন প্রমাণ দিতে হবে—প্রথমটির জন্য একটি,
দ্বিতীয়টির জন্য আরেকটি। তবে কখনো কয়েকটি
``তখনই এবং কেবল তখনই''-এর শৃঙ্খল দিয়ে $p$
থেকে $q$ পর্যন্ত যাওয়া যায়; তখন প্রতিটি ধাপ
সত্যিই দ্বিশর্তাধীন কি না নিশ্চিত করতে হবে।

\subsection{সার্বিক দাবি}

সার্বিক দাবি ব্যবহার করা সহজ: সব কিছুর জন্য যা
সত্য, নির্দিষ্ট প্রতিটির জন্যও তা সত্য। যেমন, প্রমাণের
পূর্বধারণা $A \subseteq B$ হলে $\subseteq$-র সংজ্ঞা
খুলে পাই: প্রত্যেক $x \in A$-র জন্য $x \in B$। তাই
$z \in A$ আগে জানা থাকলে $z \in B$ সিদ্ধান্ত নিতে পারি।

সার্বিক দাবি প্রমাণ করা কিছুটা কঠিন মনে হতে পারে।
সাধারণত বক্তব্যের রূপ ``$x$-এর~$P$ ধর্ম থাকলে
তার~$Q$ ধর্মও আছে'' বা ``সব $P$ হলো $Q$''।
অবশ্য সব সময় ঠিক এই রূপে থাকে না; প্রশ্নে কী
প্রমাণ করতে বলেছে, তা ধরতে কিছুটা অভ্যাস লাগে।
কিন্তু প্রায়ই একটি ধর্মবিশিষ্ট সব বস্তুর অন্য একটি
ধর্ম আছে দেখাতে হয়।

এমন দাবি প্রমাণ করতে প্রথম ধর্মবিশিষ্ট বস্তুর জন্য
একটি নাম বা চলক চালু করে দেখাই, তার দ্বিতীয় ধর্মটিও
আছে। অন্যভাবে: \emph{সব}~$P$-র জন্য প্রমাণ করতে
একটি \emph{যেকোনো}~$P$-র জন্য প্রমাণ করতে হয়।
চালু করা নামটি এমন যেকোনো~$P$-র নাম। সাধারণত
একটি অক্ষর নিই এবং প্রচলিত রীতি মানি: যেমন স্বাভাবিক
সংখ্যার জন্য $n$, !!{formula}s-এর জন্য $!A$, সেটের
জন্য $A$, অপেক্ষকের জন্য $f$ ইত্যাদি।

কৌশল হলো পুরো প্রমাণে সাধারণত্ব বজায় রাখা।
শুরুতে ধরো যে যেকোনো একটি বস্তু (``$x$'')-র
ধর্ম~$P$ আছে; তারপর শুধু সংজ্ঞা ও অনুমোদিত
পূর্বধারণা ব্যবহার করে দেখাও $x$-এর ধর্ম~$Q$-ও
আছে। $P$ ধর্ম ছাড়া $x$ কোন বিশেষ বস্তু, তা
স্থির করোনি; তাই $P$ ধর্মবিশিষ্ট সব বস্তুরই $Q$
ধর্ম আছে বলতে পারো। অর্থাৎ $x$ হলো~$P$ ধর্মবিশিষ্ট
\emph{সব} বস্তুর প্রতিনিধি।

\begin{prop}
  যেকোনো সেট $A$ ও $B$-র জন্য $A \subseteq A \cup B$।
\end{prop}

\begin{proof}
  $A$ ও $B$ যেকোনো সেট হোক। দেখাতে চাই $A
  \subseteq A \cup B$। $\subseteq$-র সংজ্ঞায় এর
  অর্থ: প্রত্যেক $x$-এর জন্য, $x \in A$ হলে
  $x \in A \cup B$। তাই $x \in A$ হোক $A$-র
  যেকোনো !!{element}। দেখাতে হবে $x \in A
  \cup B$। যেহেতু $x \in A$, তাই $x \in A$
  অথবা $x \in B$। সুতরাং $x \in
  \Setabs{x}{x \in A \lor x \in B}$। আর $\cup
  $-র সংজ্ঞায় এর অর্থ $x \in A \cup B$।
\end{proof}


\subsection{ক্ষেত্রবিচারে প্রমাণ}

ধরো, একটি বিযোজন পূর্বধারণা হিসেবে আছে, অথবা আগেই
সিদ্ধ হয়েছে: $p$ অথবা $q$ সত্য। প্রমাণ করতে চাই $r$।
দুটি ধাপে তা করো: প্রথমে $p$ সত্য ধরে~$r$ প্রমাণ,
তারপর $q$ সত্য ধরে আবার~$r$ প্রমাণ। কারণ দুটি
বিকল্পের অন্তত একটি সত্য বলে ধরা বা জানা আছে;
দুটি ধাপ দেখায়, যেকোনো একটিই~$r$-এর সত্যতার
জন্য যথেষ্ট। (দুটিই সত্য হলে $r$ সত্য হওয়ার দুটি
কারণ পাওয়া যায়। $p$ এবং~$q$ দুটিই ধরে আলাদা
করে~$r$ প্রমাণের দরকার নেই।) কী করছি বোঝাতে
``ক্ষেত্রবিচার করি'' লিখি। যেমন, জানা আছে
$x \in B \cup C$। $B \cup C$-র সংজ্ঞা
$\Setabs{x}{x \in B \text{ অথবা } x \in C}$।
অর্থাৎ সংজ্ঞা অনুযায়ী $x \in B$ অথবা $x \in C$।
এ থেকে $x \in A$ প্রমাণ করতে প্রথমে $x \in B$
ধরে তার থেকে $x \in A$ দেখাব; তারপর $x \in C$
ধরে আবার $x \in A$ দেখাব। পূর্বধারণার নিচে
``ক্ষেত্রবিচার করি'', তার নিচে ``ক্ষেত্র (1): $x \in B$'',
আর পাতার মাঝামাঝি ``ক্ষেত্র (2): $x \in C$'' লিখে
উপরের ও নিচের অর্ধেকের প্রমাণ পূরণ করতে পারো।

প্রমাণাধীন বক্তব্য নিজেও বিযোজক হলে ক্ষেত্রবিচার
বিশেষভাবে কাজে লাগে। একটি সহজ উদাহরণ:

\begin{prop}
$B \subseteq D$ এবং $C \subseteq E$ হলে
$B \cup C \subseteq D \cup E$।
\end{prop}

\begin{proof}
  (a) $B \subseteq D$ এবং (b) $C \subseteq E$
  ধরি। সংজ্ঞা অনুযায়ী প্রত্যেক $x \in B$-র জন্য
  % BN-SRC-477: restore omitted membership subjects throughout this section.
  $x \in D$ (c), আর প্রত্যেক $x \in C$-র জন্য
  $x \in E$ (d)। $B \cup C \subseteq D \cup E$
  দেখাতে $\subseteq$-র সংজ্ঞা অনুযায়ী দেখাতে হবে:
  $x \in B \cup C$ হলে $x \in D \cup E$।
  $\cup$-র সংজ্ঞায় $x \in B \cup C$ তখনই, যখন
  $x \in B$ অথবা $x \in C$। একইভাবে
  $x \in D \cup E$ তখনই, যখন $x \in D$ অথবা
  $x \in E$। সুতরাং দেখাতে হবে: যেকোনো $x$-এর
  জন্য $x \in B$ অথবা $x \in C$ হলে $x \in D$
  অথবা $x \in E$।

  \begin{quote}
  এ পর্যন্ত কেবল সংজ্ঞা খুলেছি!{} $\subseteq$ ও
  $\cup$ ছাড়া প্রস্তাবটি নতুনভাবে বলে এখন একটি
  সার্বিক শর্তাধীন বক্তব্য প্রমাণ করতে হবে। আগের
  আলোচনা অনুযায়ী, এমন একটি $x$ নিই যার জন্য
  ``যদি'' অংশটি সত্য ধরি, তারপর ``তবে'' অংশটিও
  সত্য দেখাই। অর্থাৎ $x \in B$ অথবা $x \in C$
  ধরে দেখাব $x \in D$ অথবা $x \in
  E$।\footnote{এই অনুচ্ছেদে শুধু কী করছি তার
    ব্যাখ্যা আছে—এটি প্রমাণের অংশ নয়। নিজের
    প্রমাণ লেখার সময় এত বিশদে ব্যাখ্যা দিতে হয় না।}
  \end{quote}

  ধরো $x \in B$ অথবা $x \in C$। দেখাতে হবে
  $x \in D$ অথবা $x \in E$। ক্ষেত্রবিচার করি।

  ক্ষেত্র 1: $x \in B$। (c) থেকে $x \in D$।
  তাই $x \in D$ অথবা $x \in E$। (আগের
  উপবিভাগের অনুমিতি-ধাপটিই এখানে ব্যবহার করেছি।)

  ক্ষেত্র 2: $x \in C$। (d) থেকে $x \in E$।
  তাই $x \in D$ অথবা $x \in E$।
\end{proof}


\subsection{অস্তিত্ব-দাবি প্রমাণ}

অস্তিত্ব-দাবি প্রমাণ করতে বললে প্রশ্নটি সাধারণত
``এমন একটি~$x$ আছে যে $\dots x \dots$'' আকারে
থাকে—অর্থাৎ ``$\dots x \dots$''-এ বর্ণিত ধর্মবিশিষ্ট
কোনো বস্তু আছে। তখন উপযুক্ত বস্তু চিহ্নিত করে তার
প্রয়োজনীয় ধর্মটি দেখাতে হবে। শুনতে সহজ হলেও এ ধরনের
প্রমাণ কঠিন হতে পারে। সাধারণত একটি বস্তু
\emph{নির্মাণ} বা \emph{সংজ্ঞায়িত} করে দেখাতে হয়,
এভাবে সংজ্ঞায়িত বস্তুর কাঙ্ক্ষিত ধর্ম আছে। সঠিক বস্তু
খুঁজে পাওয়া কঠিন হতে পারে, তার ধর্ম প্রমাণ করাও কঠিন
হতে পারে; কখনো আদৌ একটি বস্তু সংজ্ঞায়িত করতে
পেরেছ কি না দেখানোও কঠিন।

সাধারণত বস্তুটি নির্দিষ্ট করে লেখো, যেমন ``$x$ হোক
\dots'' (এখানে \dots{} দিয়ে কোন বস্তু বোঝাচ্ছ তা
বলবে)। প্রয়োজনে $\dots$ সত্যিই একটি বিদ্যমান বস্তু
বর্ণনা করে, তা প্রমাণ করবে; তারপর $x$-এর ধর্ম~$Q$
আছে দেখাবে। একটি সহজ উদাহরণ:

\begin{prop}
  $x \in B$ হলে এমন একটি~$A$ আছে যাতে
  $A \subseteq B$ এবং $A \neq \emptyset$।
\end{prop}

\begin{proof}
  $x \in B$ ধরি। $A = \{x\}$ নিই।
  \begin{quote}
    এখানে !!{element}s তালিকাভুক্ত করে সেট~$A$
    সংজ্ঞায়িত করেছি। $x$ একটি বস্তু বলে ধরা আছে,
    আর তালিকাভুক্ত !!{element}s দিয়ে সেট গঠন করা যায়;
    তাই এখানে সত্যিই সেট~$A$ সংজ্ঞায়িত হয়েছে কি না
    আলাদা করে দেখাতে হয় না। তবে প্রস্তাবের প্রয়োজনীয়
    ধর্মগুলি $A$-র আছে, তা এখনো দেখাতে হবে। তা ছাড়া
    প্রমাণ সম্পূর্ণ নয়!{}
  \end{quote}
  যেহেতু $x \in A$, তাই $A \neq \emptyset$।
  \begin{quote}
    এখানে $A$-কে $\{x\}$ হিসেবে সংজ্ঞায়িত করেছি,
    আর স্পষ্টতই $x \in \{x\}$ এবং
    $x \notin \emptyset$—এই তথ্যগুলি ব্যবহার করেছি।
  \end{quote}
  $x$ হলো~$\{x\}$-র একমাত্র !!{element}, এবং
  $x \in B$; তাই $A$-র প্রত্যেক !!{element}
  $B$-রও !!a{element}। $\subseteq$-র সংজ্ঞায়
  $A \subseteq B$।
\end{proof}

\subsection{অস্তিত্ব-দাবির ব্যবহার}

ধরো, একটি অস্তিত্ব-দাবি সত্য বলে জানো—প্রমাণ করেছ
অথবা ব্যবহারযোগ্য পূর্বধারণা হিসেবে পেয়েছ। যেমন,
``কোনো~$x$-এর জন্য $x \in A$'' বা ``একটি
$x \in A$ আছে।'' প্রমাণে এটি ব্যবহার করতে
চাইলে ধরে নিতে পারো, বিদ্যমান বস্তুগুলির একটির
জন্য তোমার একটি নাম আছে। $A$-তে অন্তত একটি
বস্তু আছে বলে নামটি যাকে বোঝাতে পারে, এমন
বস্তু আছে। কোনটি তা নির্দিষ্ট করে বেছে নেওয়া বা
$\in A$ ছাড়া আর কিছু বলা নাও যেতে পারে। কিন্তু
প্রমাণের জন্য একটিকে বেছে নিয়েছ বলে ধরে তাকে
নাম দিতে পারো। আগে ব্যবহার করোনি এবং
পূর্বধারণাতেও নেই, এমন নাম বেছে নেওয়া জরুরি;
নইলে ভুল হতে পারে। তাই ``কোনো $x$-এর জন্য
$x \in A$'' থেকে লেখো ``$a \in A$ হোক।''
এবার~$a$ নিয়ে অন্য পূর্বধারণাও ব্যবহার করে
একটি সিদ্ধান্ত $p$-তে পৌঁছতে পারো। $p$-তে আর
$a$ না থাকলে, $p$ নির্দিষ্ট $a \in A$ অনুমানের
উপর নির্ভর করে না; শুধু ``কোনো $x$-এর জন্য
$x \in A$'' অনুমান থেকেই তা এসেছে।

\begin{prop}
$A \neq \emptyset$ হলে $A \cup B \neq \emptyset$।
\end{prop}

\begin{proof}
  $A \neq \emptyset$ ধরি। তাই কোনো $x$-এর জন্য
  $x \in A$।
  \begin{quote}
    এখানে প্রথমে প্রস্তাবের পূর্বধারণাটিই আবার বলেছি।
    পূর্বধারণা $A \neq \emptyset$-এর মধ্যে একটি
    অস্তিত্ব-দাবি লুকিয়ে আছে; কয়েকটি সংজ্ঞা খুললে
    সেটি পাওয়া যায়। $=$-র সংজ্ঞায় $A = \emptyset$
    তখনই, যখন প্রত্যেক $x \in A$-র জন্য
    $x \in \emptyset$ এবং প্রত্যেক $x \in
    \emptyset$-র জন্য $x \in A$। দুই দিক নঞর্থক
    করলে পাই: $A \neq \emptyset$ তখনই, যখন
    হয় কোনো $x \in A$-র জন্য $x \notin \emptyset$,
    নয়তো কোনো $x \in \emptyset$-র জন্য
    $x \notin A$। $x \in \emptyset$ কখনো সত্য নয় বলে
    দ্বিতীয় বিকল্পটি অসম্ভব; আর ``$x \in A$ এবং
    $x \notin \emptyset$'' নেমে আসে শুধু
    $x \in A$-তে। তাই
    % BN-SRC-478: the nonempty set here is A, not the bound x.
    $A \neq \emptyset$ তখনই, যখন কোনো $x$-এর
    জন্য $x \in A$। এটিই অস্তিত্ব-দাবি। এবার
    % BN-SRC-834: introduce one witness from the plural collection.
    $A$-র !!{element}s-এর একটির নাম চালু করে
    দাবিটি ব্যবহার করি:
  \end{quote}
  $a \in A$ হোক।
  \begin{quote}
    $\in A$ ধর্মবিশিষ্ট একটি বস্তুর নাম দিয়েছি।
    এবার~$a$ নিয়ে যুক্তি চালাব, কিন্তু শুধু
    $a \in A$-ই ধরব; অতিরিক্ত কিছু নয়:
  \end{quote}
  $a \in A$ বলে $\cup$-র সংজ্ঞায়
  $a \in A \cup B$। তাই কোনো $x$-এর জন্য
  $x \in A \cup B$; অর্থাৎ $A \cup B \neq \emptyset$।
  \begin{quote}
    শেষ ধাপে ``$a \in A \cup B$'' থেকে
    ``কোনো $x$-এর জন্য $x \in A \cup B$'' বলেছি।
    শেষের কথাটিতে আর $a$ নেই; তাই ``কোনো $x$-এর
    জন্য $x \in A$'' থেকেই ``কোনো $x$-এর জন্য
    $x \in A \cup B$'' আসে। আর এর অর্থ
    $A \cup B \neq \emptyset$।
  \end{quote}
\end{proof}

উপরের মতো বদ্ধ চলক ``$x$'' এবং অনুমান থেকে
নেওয়া নাম $a$ আলাদা রাখা ভালো অভ্যাস। তবে
বাস্তবে প্রায়ই তা করি না; শুধু $x$ লিখি:
\begin{quote}
$A \neq \emptyset$ ধরি; অর্থাৎ একটি $x \in A$
আছে। $\cup$-র সংজ্ঞায় $x \in A \cup B$।
তাই $A \cup B \neq \emptyset$।
\end{quote}
এভাবে লিখলে কিন্তু আলাদা অস্তিত্ব-দাবির জন্য আলাদা
$x$, $y$ ব্যবহারে বিশেষ সতর্ক হতে হবে। যেমন,
``$A \neq \emptyset$ এবং $B \neq \emptyset$ হলে
$A \cap B \neq \emptyset$''—এই মিথ্যা দাবির
নিচের ``প্রমাণ''টি \emph{ঠিক নয়}:
\begin{quote}
$A \neq \emptyset$ এবং $B \neq \emptyset$ ধরি।
তাই কোনো $x$-এর জন্য $x \in A$, আবার কোনো
$x$-এর জন্য $x \in B$। যেহেতু $x \in A$
এবং $x \in B$, তাই $\cap$-র সংজ্ঞায়
$x \in A \cap B$। অতএব $A \cap B \neq
\emptyset$।
\end{quote}
ভুল ধাপটি কোথায় এবং সিদ্ধান্তটি কেন আসে না,
তা কি বলতে পারো?

\end{document}
